\documentclass[a4paper,11pt]{article}

% --- Pakker for matematikk og font ---
\usepackage[english]{babel}
\usepackage{amsmath, amsthm, amssymb}
\usepackage{mathtools} 
\usepackage{geometry} 
\usepackage{enumitem}
\geometry{margin=2cm}

\setlength{\parindent}{0pt}
\setlength{\parskip}{0.5em} 

% make amsthm use a filled in black square
\renewcommand{\qedsymbol}{$\blacksquare$}

% Mengder
\newcommand{\N}{\mathbb{N}}
\newcommand{\Z}{\mathbb{Z}}
\newcommand{\Q}{\mathbb{Q}}
\newcommand{\R}{\mathbb{R}}
\newcommand{\C}{\mathbb{C}}

% Topologi/Analyse
\newcommand{\inv}[1]{#1^{-1}} 
\newcommand{\interior}[1]{#1^\circ} 
\newcommand{\closure}[1]{\overline{#1}} 
\newcommand{\boundary}{\partial} 

% Greske bokstaver
\newcommand{\eps}{\varepsilon}

% --- Document Info ---
\title{MAT2410 - Introduction to complex analysis: \\ Mandatory assignment 2 of 2}
\author{Vegard Nyquist Borgeraas}
\date{Autumn 2026}

\begin{document}

\maketitle

\section*{Problem 1}
Let $f(z)=\frac{e^{z}}{(1-z)^{2}}$, and let $\sum_{k=0}^{\infty}a_{k}z^{k}$ be the power series expansion of $f$ around $0$.

\begin{enumerate}[label=\textbf{(\alph*)}]
    \item Determine $a_{0}$, $a_{1}$ and $a_{2}$.
    \begin{proof}
        Using the formula $a_k = \frac{f^{(k)}(0)}{k!}$, we have:
        \begin{align*}
            a_{0} &= f(0) = \left. \frac{e^z}{(1-z)^2} \right|_{z=0} = 1 \\
            a_{1} &= f'(0) = \left. \frac{e^z(3-z)}{(1-z)^3} \right|_{z=0} = 3 \\
            a_{2} &= \frac{f''(0)}{2} = \frac{1}{2} \left. \frac{e^z(z^2-6z+11)}{(1-z)^4} \right|_{z=0} = \frac{11}{2}
        \end{align*}
    \end{proof}

    \item Determine the radius of convergence of the power series.
    \begin{proof}
        The radius of convergence of the power series at $z=0$ is equal to the distance to the nearest singularity of $f$. Since the only singularity is a pole at $z=1$, the radius of convergence is $1$.
    \end{proof}
    
    \item Compute $\int_{|z|=2}f(z)\,dz$, where the path around the circle $|z|=2$ is oriented counterclockwise.
    \begin{proof}
        Let $g(z) = e^z$. By applying Cauchy's integral formula for derivatives we get:
        \[
            e = g'(1) = \frac{1}{2\pi i}\oint_{|z|=2} \frac{g(z)}{(z-1)^2}\,dz = \frac{1}{2\pi i}\oint_{|z|=2} \frac{e^z}{(1-z)^2}\,dz.
        \]
        Rearranging this equation, we get:
        \[
            \oint_{|z|=2} f(z)\,dz = 2\pi i e.
        \]
    \end{proof}
\end{enumerate}

\section*{Problem 2}
Let $D\subseteq\C$ be a bounded domain with a piecewise smooth boundary. Let $f$ be analytic on $D$, extending smoothly to $\closure{D}$.

\begin{enumerate}[label=\textbf{(\alph*)}]
    \item Show that if $f$ is constant on $\boundary D$, then $f$ is constant on $\closure{D}$.
    \begin{proof}
        Since $f$ is constant on $\boundary D$, there exists a constant $c\in\C$ such that $f(z) = c$ for all $z\in\boundary D$.
        Let $z \in D$. By the Cauchy integral formula, we have:
        \[
        f(z) = \frac{1}{2\pi i} \oint_{\boundary D} \frac{f(w)}{w - z}\,dw = \frac{c}{2\pi i} \oint_{\boundary D} \frac{1}{w - z}\,dw = c
        \]
        Thus, $f(z) = c$ for all $z \in D$. Since $f$ equals the same constant on both $D$ and $\boundary D$, it is constant on $\closure{D}$.
    \end{proof}
    
    \item Show that if $f(z)\ne0$ for all $z\in\boundary D$ and
    \[ \int_{\boundary D}\frac{1}{f(z)}\,dz\neq0, \]
    then $f$ has a zero in $D$.
    \begin{proof}
        We prove the contrapositive.
        Assume that $f$ has no zeros in $D$. 
        Since $f$ is also non-zero on $\boundary D$, the function $\frac{1}{f(z)}$ is analytic on $D$ and extends smoothly to $\closure{D}$.
        By Cauchy's integral theorem, we get: 
        \[ \int_{\boundary D}\frac{1}{f(z)}\,dz = 0. \]
    \end{proof}
\end{enumerate}

\section*{Problem 3}
Let $\alpha\in\C$. For each integer $k\ge1$, let
\[ \binom{\alpha}{k}=\frac{\alpha(\alpha-1)\dots(\alpha-k+1)}{k!}, \]
and let $\binom{\alpha}{0}=1$. Consider the power series
\[ f(z)=\sum_{k=0}^{\infty}\binom{\alpha}{k}z^{k}. \]

\begin{enumerate}[label=\textbf{(\alph*)}]
    \item Determine the radius of convergence of the power series; the answer will depend on whether $\alpha\in\Z_{\ge0}$ or not.
    \begin{proof}
        If $\alpha\in\Z_{\ge0}$, we have $\binom{\alpha}{k} = 0$ for all $k>\alpha$. In this case, the power series terminates and reduces to a polynomial, which has a radius of convergence of $\infty$. 
        
        Otherwise, if $\alpha \in \C \setminus \Z_{\ge0}$, then $\binom{\alpha}{k} \neq 0$ for all $k \ge 0$. Applying the ratio test, we get:
        \[
            \lim_{k \to \infty} \left| \frac{\binom{\alpha}{k}}{\binom{\alpha}{k+1}} \right| 
            = \lim_{k \to \infty} \left| \frac{k+1}{\alpha - k} \right| = 1.
        \]
        Therefore, the radius of convergence is $1$ for $\alpha \notin \Z_{\ge0}$.
    \end{proof}

    
\item Recall that for $w\in\C\setminus\R_{\le0}$ we define the principal branch $P_{\alpha}$ of the $\alpha$-th power function $w^{\alpha}$ by 
    \[ P_{\alpha}(w)=e^{\alpha \operatorname{Log}(w)}, \]
    where $\operatorname{Log}(w)$ denotes the principal branch of the logarithm. Show that
    \[ P_{\alpha}'(w)=\alpha P_{\alpha-1}(w) \]
    for all $w\in\C\setminus\R_{\le0}$.
    \begin{proof}
        Let $w\in\C\setminus\R_{\le0}$. Applying the chain rule, we have:
        \begin{align*}
            P_{\alpha}'(w) &= \frac{d}{dw} e^{\alpha \operatorname{Log}(w)} \\
            &= e^{\alpha \operatorname{Log}(w)} \cdot \frac{\alpha}{w} \\
            &= \alpha e^{\alpha \operatorname{Log}(w)} e^{-\operatorname{Log}(w)} \\
            &= \alpha e^{(\alpha-1) \operatorname{Log}(w)} \\
            &= \alpha P_{\alpha-1}(w).
        \end{align*}
    \end{proof}
    
    \item Determine $P_{\alpha}^{(k)}(1)$ for all $k \ge 0$ and use this to show that for $|z|<1$, we have
    \[ f(z)=P_{\alpha}(1+z). \]
    \begin{proof}
        By the recursive formula derived in the previous part, it follows that $P_{\alpha}^{(k)}(1) = \alpha(\alpha - 1) \cdots (\alpha - k + 1)$ for $k \ge 1$, and define $P_{\alpha}^{(0)}(1) = 1$.
        Thus, the power series expansion of $P_{\alpha}(z+1)$ around $z = 0$ for $|z|<1$ gives:
        \[ 
            P_{\alpha}(z+1)
            = \sum_{k=0}^{\infty} \frac{P_{\alpha}^{(k)}(1)}{k!} z^k 
            = \sum_{k=0}^{\infty} \frac{\alpha(\alpha - 1) \cdots (\alpha - k + 1)}{k!} z^k
            = \sum_{k=0}^{\infty} \binom{\alpha}{k} z^k
            = f(z).
        \]
    \end{proof}
\end{enumerate}

\section*{Problem 4}
Let $f(z)=\frac{e^{z}}{z}$, let $r>0$, and let $\gamma_{r}$ be the path
\[ \gamma_{r}(t)=-r\sin(t)+ir\cos(t), \quad 0\le t\le\pi. \]
Show that
\[ \left|\int_{\gamma_{r}}f(z)\,dz\right|\le\pi. \]
\begin{proof}
    The length of the path is $L:=\pi r$. Let $z$ be in the image of $\gamma_r$. We have:
    \[ 
        |f(z)| = \left|\frac{e^z}{z}\right| = \frac{|e^z|}{|z|} = \frac{e^{\text{Re}(z)}}{r} \leq \frac{e^0}{r} = \frac{1}{r} =: M.
    \]
    By the ML-estimate, it follows that:
    \[ \left|\int_{\gamma_{r}}f(z)\,dz\right| \le M L = \pi. \]
\end{proof}

\section*{Problem 5}
For each of the hypothetical function descriptions below, either give an example of a function satisfying the description, or prove that no such function can exist.

\begin{enumerate}[label=\textbf{(\alph*)}]
    \item A function which is defined and analytic on $\C\setminus\Z$, but does not admit an analytic extension to any point of $\Z$.
    \begin{proof}
        Let $f(z) = \frac{1}{\sin(\pi z)}$. This function has non-removable poles exactly at the integers, and therefore cannot be analytically extended to any point in $\Z$.
    \end{proof}
    
    \item An entire function such that $|f^{(k)}(0)|\ge k!$ for all $k\ge0$.
    \begin{proof}
        Suppose for the sake of contradiction that there exists such a function. Consider its power series expansion around $z=0$. Since $f$ is entire, the radius of convergence is $\infty$:
        \[\sum_{k=0}^{\infty} \frac{f^{(k)}(0)}{k!}z^k\]
        This series must converge absolutely for all $z\in\C$, and in particular for $z=1$. However, we have:
        \[
        \sum_{k=0}^{\infty} \left|\frac{f^{(k)}(0)}{k!}1^k\right| 
        \geq \sum_{k=0}^{\infty} \frac{k!}{k!} 
        = \sum_{k=0}^{\infty} 1 = \infty,
        \]
        which diverges. This is a contradiction.
    \end{proof}
    
    \item An entire function such that $f\left(\frac{1}{n}\right)=\frac{1}{2-1/n}$ for all integers $n\ge1$.
        \begin{proof}
        Suppose for contradiction that there exists such an entire function $f$.
        Define $g(z) = \frac{1}{2-z}$. Then $g\left(\frac{1}{n}\right) = \frac{1}{2-1/n}$ for all $n\geq1$.
        Since $f$ is continuous, we have: 
        \[ f(0) = \lim_{n\to\infty} f\left(\frac{1}{n}\right) = \frac{1}{2} = g(0). \]
        Since every neighborhood of $0$ contains points of the sequence $(1/n)$, the set of points where $f$ and $g$ agree contains a non-isolated point at $z=0$.
        By the Uniqueness Principle $f=g$. 
        However, this is a contradiction since $g$ has a pole at $z=2$ and therefore cannot be extended to $\C$.
    \end{proof}
\end{enumerate}

\end{document}